(1) 均匀细梁长度 \(L\)、截面积 \(A\)、截面二次矩 \(I\)、密度 \(\rho\)、杨氏模量 \(E\)。Rayleigh梁保留截面转动惯量但忽略剪切变形。由动能 \(T=\frac12\int_0^L(\rho A u_t^2+\rho I u_{xt}^2)dx\) 与弯曲势能 \(V=\frac12\int_0^LEIu_{xx}^2dx\) 推导横向位移 \(u(x,t)\) 的场方程。
(2) 在无限梁上代入 \(u\propto e^{i(kx-\omega t)}\) 求色散 \(\omega(k)\)，分别给出长波弯曲波和短波的相、群速度渐近。
(3) 对简支端条件 \(u=u_{xx}=0\) 于 \(x=0,L\)，求全部本征频率并比较忽略转动惯量的Euler–Bernoulli预测，证明Rayleigh频率逐模不高于后者。
(4) 推导无外力无阻尼时机械能守恒的边界通量表达式，说明为什么所有非平移本征模的 \(\omega_n^2>0\)，并指出极短波时剪切变形忽略不再可靠。
